\section{Introduction}%
\label{sec:intro}

We are interested in estimating point correspondences in videos containing dynamic objects and a moving camera.
Establishing point correspondences is an important problem in computer vision with many applications~\cite{agarwal2011building, teed2021droid, jiang2018super, yang2022banmo}.
There are two variants of this problem.
In \emph{optical flow}, the goal is to estimate the velocities of all points within a video frame.
This estimation is performed jointly for all points, but the motion is only predicted at an infinitesimal distance.
In \emph{point tracking}, the goal is to estimate the motion of points over an extended period of time.
For efficiency, trackers usually focus on a sparse set of points and treat them as statistically independent.
This is the case even for recent techniques such as TAPIR~\cite{doersch2023tapir} and PIPs++~\cite{zheng2023pointodyssey}, which employ modern architectures such as transformers and can track points even in the presence of occlusions.
However, points often have strong statistical dependencies (\eg, because they belong to the same object), which offers an opportunity for improvement.

We hypothesise that accounting for the dependency between tracked points can significantly improve tracking performance.
To explore this idea, we introduce \method, a new tracker that supports joint estimation of a very large number of tracks simultaneously, utilizing a transformer-based architecture and attention between tracks.
As we show in ablations, tracking points jointly significantly improves the tracking accuracy, especially when points are occluded.
In fact, we find that we can improve tracking quality further by adding \emph{more} points to the tracker than required by the user.
These additional \emph{support points} expand the context of the tracker, porting to point tracking the idea of using context, common in visual object tracking~\cite{jia2012visual, li2015reliable, song2017crest}.

In addition to joint tracking, \method incorporates additional architectural design innovations.
The network is a transformer operating in a sliding window fashion on a two-dimensional token representation, where dimensions are time and the set of tracked points~\cite{gberta_2021_ICML}.
Via suitable self-attention operators, the transformer considers each track as a whole for the duration of a window and can exchange information between tracks, thus exploiting their dependencies.
However, attention can be expensive when there are many tracks.
Hence, we introduce the concept of proxy tokens to point tracking that act similarly to registers~\cite{darcet2023vision, jaegle22perceiver, nagrani2021attention} while also reducing the memory complexity.
These tokens are processed as if they were a small number of additional tracks and allow the switch from expensive self-attention between tracks to efficient cross-attention between tracks and proxies.
In this way, \method can jointly track a \emph{near-dense set of tracks} on a single GPU\@ at inference.

\method is designed as an online tracker that operates on relatively short windows of frames.
Within a window, the tracks are initialized with queried points, and the network is tasked with progressively refining these initial estimations through iterative applications of the transformer.
The windows partially overlap and communicate, similar to a recurrent network.
Each subsequent overlapping window starts with refined predictions of the previous window and updates tracks for the new frames.
We optimize the recurrent application of the network via unrolled training, porting this concept from recurrent networks and visual object tracking~\cite{held2016learning, girshick2015deformable} to point tracking.
In this way,
\method achieves \emph{excellent long-term tracking performance} as well; in particular, joint and recurrent tracking allows points to be tracked through long occlusions.

We train \method on the synthetic TAP-Vid-Kubric~\cite{tapvid} and evaluate it on  TAP-Vid-\{DAVIS, RGB-Stacking\}, PointOdyssey~\cite{zheng2023pointodyssey} and DynamicReplica~\cite{karaev2023dynamicstereo}.
Our architecture works well for tracking single points and excels for groups of points, obtaining state-of-the-art tracking performance in several benchmarks compared to prior trackers.
In summary, our contributions are the following:
(i) we introduce the concept of \emph{joint point tracking} by sharing information between tracked points through an attention mechanism;
(ii) we propose to use \emph{support points} to provide additional context for point tracking;
(iii) we \emph{reduce the model's memory complexity} with proxy tokens; and
(iv) we propose an \emph{unrolled model training strategy}, which further improves tracking and occlusion accuracy.

\begin{figure}[t]
\includegraphics[width=\textwidth]{figures/plot_joint_individual.pdf}
\vspace{-2em}
\caption{\textbf{Non-joint tracking (top row) vs.~joint tracking (bottom).}  Background points are colored cyan while foreground points are magenta.
Non-joint point tracking causes background points to follow object  motion (cyan points in later frames), while tracking points together results in better tracks for foreground and background points.}%
\label{fig:single_joint}
\end{figure}
